% !PRINTED 771
\weberpage{771}

Um die Irreducibilität für ein beliebig zusammengesetztes $n$
nachzuweisen, wenden wir die vollständige Induction an. Wir
setzen
\[
n = p^k n',
\]
und verstehen unter $p$ eine Primzahl, die in $n$, aber nicht in $n'$
aufgeht, und setzen die \gesperrt{Irreducibilität von} $f_{n'}(x)$ \gesperrt{schon als
bewiesen voraus}. Ist dann $f_n(x)$ in zwei rationale Factoren
$\varphi(x)$, $\psi(x)$ zerlegbar, deren keiner constant ist:
\weq{7}{f_n(x) = \varphi(x)\, \psi(x),}
so leiten wir aus $\varphi(x)$ die Function $\Phi(x)$ ab, deren Wurzeln
die $p^{k\,\mathrm{ten}}$ Potenzen der Wurzeln von $\varphi(x)$ sind, die der Bedingung
\[
\varphi(x) \equiv \Phi(x) \ (\mathrm{mod}\ p)
\]
genügt, oder auch, indem wir mit $\chi(x)$ eine zweite ganzzahlige
Function von $x$ bezeichnen:
\[
\varphi(x) = \Phi(x) + p\, \chi(x).
\]

Die Wurzeln von $\Phi(x)$ sind aber primitive Einheitswurzeln
vom Grade $n'$, und daher muss unter diesen wenigstens eine
sein, $\alpha'$, für die $\Phi(\alpha')$ verschwindet, also
\[
\varphi(\alpha') = p\, \chi(\alpha')
\]
ist. Wegen der vorausgesetzten Irreducibilität von $f_{n'}(x)$ muss
aber diese Gleichung für \gesperrt{alle} Einheitswurzeln $\alpha'$ bestehen, und
aus dem gleichen Grunde ist, wenn $\vartheta(x)$ eine ganzzahlige
Function ist, für alle $\alpha'$:
\[
\psi(\alpha') = p\, \vartheta(\alpha').
\]

Hiernach folgt aus (7):
\[
f_n(\alpha') = p^2\, \chi(\alpha')\, \vartheta(\alpha'),
\]
und nach (4):
\[
1 = p\, \chi(\alpha')\, \vartheta(\alpha')\, g(\alpha').
\]

Das Product $\chi(x)\, \vartheta(x)\, g(x)$ ist aber eine ganze ganzzahlige
Function von $x$, deren Rest bei der Division mit $f_{n'}(x)$ gleichfalls ganzzahlige Coëfficienten hat und mit $\Theta(x)$ bezeichnet sein
mag. Dieser Rest ist von niedrigerem Grade, als $f_{n'}(x)$, und die
Function
\[
p\, \Theta(x) - 1
\]
verschwindet für \gesperrt{alle Wurzeln} von $f_{n'}(x)$. Folglich muss die
Gleichung
\[
p\, \Theta(x) = 1
\]
identisch sein, was aber offenbar unmöglich ist, da $\Theta(x)$ für ein
%CONT>>
